\documentclass[7pt,a4paper]{article}
\usepackage[a4paper,landscape,margin=0.4cm,top=0.3cm,bottom=0.3cm]{geometry}
\usepackage[T1]{fontenc}
\usepackage[utf8]{inputenc}
\usepackage{lmodern}
\usepackage{amsmath,amssymb}
\usepackage{multicol}
\usepackage{xcolor}
\setlength{\parindent}{0pt}
\setlength{\parskip}{0pt}
\setlength{\baselineskip}{8pt}
\pagestyle{empty}
\linespread{0.9}

\newcommand{\E}{\mathbf{E}}
\newcommand{\B}{\mathbf{B}}
\newcommand{\D}{\mathbf{D}}
\renewcommand{\H}{\mathbf{H}}

\begin{document}
\scriptsize

% COVER PAGE WITH EXAM GUIDELINES
\begin{multicols*}{2}

{\Large \textbf{REM Exam Sheet -- Examen Referentie}}\\
{\small \textit{Gebasseerd op officiële examen-richtlijnen}}\\[3pt]

\textbf{10 Mandatory Proof Topics (LEREN!):}\\
\fbox{\begin{minipage}{0.48\textwidth}
\tiny
1. $\nabla \times \E = 0$ admits electric potential (H2.3.1)\\
2. Ampère-Maxwell consistency + continuity (H7-H8)\\
3. Poynting's theorem proof (H8.1.2)\\
4. Gauge transformations preserve fields (H10.1.2)\\
5. Retarded potentials satisfy wave eq. (H10.2.1)\\
6. Liénard-Wiechert: steps + interpretation\\
7. Time dilation, length contraction, velocity add. (H12)\\
8. Covariant Maxwell from tensor notation (H12)\\
9. Bound charges from polarization (H4.2.1)\\
10. Dipole radiation + approximations (H11.1.2)
\end{minipage}}\\[3pt]

\columnbreak

\textbf{Examen-Gestuurde Secties:}\\
\textbf{\textcolor{red}{ZEER BELANGRIJK (100\%):}} H2, H5, H12\\
\textbf{\textcolor{orange}{BELANGRIJK (80\%):}} H3, H4, H6, H7, H8, H10\\
\textbf{\textcolor{blue}{MEDIUM (60\%):}} H9, H11\\
\\
\textbf{SKIP:} Capacitance (2.5.4), Multipole (5.4.3), \\
Inductance (7.2.3), Stress tensor (8.2), \\
Ferromagnetism (6.4.2)\\

\end{multicols*}

\vspace*{-3pt}
\hrule
\vspace*{2pt}

% PAGE 1: H2-H6 DETAILED CONTENT
\begin{multicols*}{5}


\textbf{H2: Electrostatica}\\\textcolor{gray}{{Sect: Entire chapter}}\\\vspace{0.5pt}
$\bullet$ H2: Electrostatics 
2.1 The electric field 
2.1.1 introduction\\
$\bullet$ > the test charge may be moving 
2.1.2 Coulombs law 
Coulombs law\\
$\bullet$ 2.1.2 Coulombs law 
Coulombs law\\
$\bullet$ with $\epsilon$0 the permittivity of free space and r:    the relative position 
2.1.3 the electric field 
...\\
$\bullet$ electric field = field only dependent on source charges 
> tells us how a test charge of Q=1 will move thro...\\
$\bullet$ defined via: 
properties of the electric field 1: E is a vector field dependent on r the position in space ...\\
$\bullet$ -     >> Coulomb's law\\
$\bullet$ 2.2 divergence and curl of electrostatic fields 
2.2.1 field lines, flux and Gauss's law\\
$\bullet$ 2.2 divergence and curl of electrostatic fields 
2.2.1 field lines, flux and Gauss's law 
flux of E the flu...\\
$\bullet$ Gauss's law for one charge at the origin we have, through a sphere of radius r:\\
\columnbreak

\textbf{H3: Potentialen}\\\textcolor{gray}{{Sect: 3.2, 3.3.2, 3.4}}\\\vspace{0.5pt}
$\bullet$ H3: Potentials 
3.2 the method of images 
charge above a conductor Consider a point charge q held a distanc...\\
$\bullet$ > the potential above the plate is determined by q, but also the induced charge on the plate  
 > this is a...\\
$\bullet$ with boundary conditions\\
$\bullet$ now we look at the boundary conditions 
> (iv) requires A=0\\
$\bullet$ they will still satisfy Laplace's equation:\\
$\bullet$ to determine Cn we use the (iii) boundary condition 
> this requires:\\
$\bullet$ symmetry  
Consider the Laplace equation in spherical coordinates:\\
$\bullet$ 3.4.1 approximate potentials at large distances 
electric dipole consider two opposite charges +-q separate...\\
\columnbreak

\textbf{H4: E-velden in materie}\\\textcolor{gray}{{Sect: 4.1.4, 4.2, 4.3, 4.4.1, 4.4.2}}\\\vspace{0.5pt}
$\bullet$ H4: Electric fields in matter 
4.1 polarization 
4.1.1 dielectrics\\
$\bullet$ 4.1.4 polarization 
polarization P place a material in a uniform field\\
$\bullet$ 4.1.4 polarization 
polarization P place a material in a uniform field 
> if the material is neutral, there...\\
$\bullet$ > either way, there will be a lot of dipoles pointing along the direction of the field  
  > the material b...\\
$\bullet$ 4.2 the field of a polarized object 
4.2.1 bound charges 
potential of a polarized object Consider a polari...\\
$\bullet$ 4.2.2 physical interpretation of bound charges 
interpretation of bound\\
$\bullet$ >> the net charge is the bound charge, because it cannot be removed 
    > in a dielectric, every e- is bou...\\
$\bullet$ > in a dielectric, every e- is bound to an atom 
bound charge mathematically to calculate the amount of bou...\\
$\bullet$ > we can rewrite this as qd 
  > the bound charge that piles up at the righthand of the tube is therefore:\\
$\bullet$ >> the effect of polarization is to paint a bound charge $\sigma$=Pn over the surface\\
\columnbreak

\textbf{H5: Magnetostatica}\\\textcolor{gray}{{Sect: Entire chapter}}\\\vspace{0.5pt}
$\bullet$ H5: Magnetostatics 
5.1 the Lorentz force law 
5.1.2 magnetic forces\\
$\bullet$ 5.1.2 magnetic forces 
Lorentz force law The magnetic force on a charge Q with velocity v in a magnetic fie...\\
$\bullet$ 5.2 the Biot-Savart law 
5.2.1 steady currents\\
$\bullet$ 5.2.2 the magnetic field of a steady current 
Biot-Savart law\\
$\bullet$ units: Tesla: 1T = 1N/(A.m) 
Surface and volume Biot-Savart for a surface:\\
$\bullet$ 5.3.2 the divergence and curl of B 
divergence of B  Biot-Savart:\\
$\bullet$ divergence of B apply the curl to Biot-Savart:\\
$\bullet$ but we are integrating over the volume that appears in the Biot-Savart law 
> this is large enough to fit a...\\
$\bullet$ we also have: F = Q(E + vxB) 
5.4 the magnetic vector potential 
5.4.1 the vector potential\\
$\bullet$ 5.4 the magnetic vector potential 
5.4.1 the vector potential 
vector potential A Since $\nabla$B=0, we can...\\
\columnbreak

\textbf{H6: B-velden in materie}\\\textcolor{gray}{{Sect: 6.1.4, 6.2, 6.3, 6.4.1}}\\\vspace{0.5pt}
$\bullet$ H6: Magnetic fields in matter 
6.1 magnetization 
6.1.1 diamagnets, paramagnets, ferromagnets\\
$\bullet$ > these are magnetic dipoles 
magnetization normally the atomic magnetic dipoles are randomly orientated  
...\\
$\bullet$ > medium becomes magnetically polarized 
types of magnetization if we apply a magnetic field B, the we have...\\
$\bullet$ - diamagnets: material magnetizes opposite to B 
 - ferromagnet: material retains magnetization after B is ...\\
$\bullet$ >> the magnetization isn't only dependent on the current applied field B 
    > it is also dependent on the...\\
$\bullet$ 6.1.4 magnetization 
magnetization M = magnetic dipole moment per unit volume\\
$\bullet$ 6.1.4 magnetization 
magnetization M = magnetic dipole moment per unit volume 
6.2 the field of a magnetize...\\
$\bullet$ 6.2 the field of a magnetized object 
6.2.1 bound currents 
potential of magnetized\\
$\bullet$ object 
consider a piece of magnetized material with magnetization M 
> the vector potential of a single di...\\
$\bullet$ define the bound currents: 
    and\\
\columnbreak


\end{multicols*}

\newpage

% PAGE 2: H7-H12 DETAILED CONTENT  
\begin{multicols*}{5}


\textbf{H7: Elektrodynamica}\\\textcolor{gray}{{Sect: 7.1.2, 7.1.3, 7.2, 7.3}}\\\vspace{0.5pt}
$\bullet$ 7.2 electromagnetic induction 
7.2.1 Faraday's law 
Faraday's law a changing magnetic field induces an elec...\\
$\bullet$ 7.2.1 Faraday's law 
Faraday's law a changing magnetic field induces an electric field 
> now we have, for ...\\
$\bullet$ 7.2.2 the induced electric field 
Biot-Savart for E If we have E in a pure Faraday field, ie E purely due t...\\
$\bullet$ 7.3 Maxwell's equations 
7.3.1 electrodynamics before Maxwell\\
$\bullet$ 7.3 Maxwell's equations 
7.3.1 electrodynamics before Maxwell 
problem with Ampère's law for the equation:\\
$\bullet$ > Ienc = 0 which we cant find on the lefthand side 
7.3.2 how Maxwell fixed Ampère's law 
Maxwell's equatio...\\
$\bullet$ 7.3.2 how Maxwell fixed Ampère's law 
Maxwell's equation We can make the righthand side of the equation zer...\\
$\bullet$ 7.3.3 Maxwell's equations 
Maxwell's equations\\
$\bullet$ 7.3.3 Maxwell's equations 
Maxwell's equations\\
$\bullet$ continuity equation\\
\columnbreak

\textbf{H8: Conservatiewetten}\\\textcolor{gray}{{Sect: 8.1, 8.2.1, 8.2.2, 8.2.3}}\\\vspace{0.5pt}
$\bullet$ 8.1 charge and energy 
8.1.1 the continuity equation 
continuity equation The charge in a volume V is:\\
$\bullet$ 8.1.1 the continuity equation 
continuity equation The charge in a volume V is:\\
$\bullet$ 8.1.2 Poynting's theorem 
Energy in E and B The work needed to assemble a static charge distribution is giv...\\
$\bullet$ Poynting's theorem If a configuration of charges changes a bit over a time dt, then we have: 
  dW =\\
$\bullet$ energy remaining in the fields, less the energy that flowed out through the surface 
Poynting vector Define...\\
$\bullet$ then Poynting's theorem becomes:\\
$\bullet$ continuity equation for energy if there is no work done, thus dW/dt=0, we find that:\\
$\bullet$ 8.2 momentum 
8.2.1 newton's third law in electrodynamics\\
\columnbreak

\textbf{H9: EM-golven}\\\textcolor{gray}{{Sect: 9.2.3, 9.3.1 only}}\\\vspace{0.5pt}
$\bullet$ 9.2 electromagnetic waves in a vacuum 
9.2.1 the wave equation for E and B 
wave equation for E and B In re...\\
$\bullet$ 9.2.1 the wave equation for E and B 
wave equation for E and B In regions of space with no charge or curren...\\
$\bullet$ which are two wave equations of waves travelling at the speed:\\
$\bullet$ 9.2.2 monochromatic plane waves 
plane waves in E and B suppose we have two plane waves travelling in z-dir...\\
$\bullet$ 9.2.2 monochromatic plane waves 
plane waves in E and B suppose we have two plane waves travelling in z-dir...\\
$\bullet$ >> the E and B are transverse waves, 
        ie: their amplitudes are perpendicular to the direction of pr...\\
$\bullet$ with k the vector of magnitude k pointing in the direction of propagation 
         n the polarization vector\\
$\bullet$ waves 
If we have plane waves for E and B, we also know: 
     with\\
$\bullet$ 9.3 electromagnetic waves in matter 
9.3.1 propagation in linear media 
Maxwell eqn in linear media inside ...\\
$\bullet$ energy density,  
Poynting vector and intensity 
if we just replace µ0>µ and $\epsilon$0>$\epsilon$ then:\\
\columnbreak

\textbf{H10: Potentialen & velden}\\\textcolor{gray}{{Sect: 10.1, 10.2, 10.3}}\\\vspace{0.5pt}
$\bullet$ 10.1.2 Gauge transformations 
Gauge transformations suppose we have two potentials (V,A) and (V',A') for th...\\
$\bullet$ 10.1.2 Gauge transformations 
Gauge transformations suppose we have two potentials (V,A) and (V',A') for th...\\
$\bullet$ where we can absorb k(t) inside $\lambda$, which wont affect the gradient of $\lambda$ 
> we therefore find...\\
$\bullet$ 10.1.3 Coulomb gauge and Lorentz gauge 
Coulomb gauge The defining choice for this gauge is:\\
$\bullet$ 10.1.3 Coulomb gauge and Lorentz gauge 
Coulomb gauge The defining choice for this gauge is:\\
$\bullet$ > -$\nabla$V-$\partial$A/$\partial$t doesn't change instantly 
Lorenz gauge the defining choice for this ga...\\
$\bullet$ >> this forms an inhomogeneous wave equation 
10.1.4 Lorentz force law in potential form\\
$\bullet$ 10.2 continuous distributions 
10.2.1 retarded potentials 
retarded time For a moving charge, its not the s...\\
$\bullet$ 10.2.1 retarded potentials 
retarded time For a moving charge, its not the status of the source right now t...\\
$\bullet$ retarded potentials We have to account for the retarded time in nonstatic sources:\\
\columnbreak

\textbf{H11: Straling}\\\textcolor{gray}{{Sect: 11.1.1, 11.1.2}}\\\vspace{0.5pt}
$\bullet$ H11: Radiation 
11.1 dipole radiation\\
$\bullet$ H11: Radiation 
11.1 dipole radiation 
11.1.1 what is radiation?\\
$\bullet$ 11.1 dipole radiation 
11.1.1 what is radiation? 
radiation when charges accelerate their fields can transp...\\
$\bullet$ 11.1.1 what is radiation? 
radiation when charges accelerate their fields can transport energy irreversibly...\\
$\bullet$ > imagine a giant sphere of radius r around this source 
 > the power passing through its surface is the in...\\
$\bullet$ radiation of static sources The area of a sphere is 4$\pi$r$^2$ 
> for radiation to occur the Poynting vect...\\
$\bullet$ radiation of static sources The area of a sphere is 4$\pi$r$^2$ 
> for radiation to occur the Poynting vect...\\
$\bullet$ >> for radiation we need to consider E and B fields that go like 1/r 
11.1.2 electric dipole radiation\\
$\bullet$ >> for radiation we need to consider E and B fields that go like 1/r 
11.1.2 electric dipole radiation 
osc...\\
$\bullet$ 11.1.2 electric dipole radiation 
oscillating dipole consider a dipole spinning around its axis at a angula...\\
\columnbreak

\textbf{H12: Relativiteit}\\\textcolor{gray}{{Sect: Entire chapter}}\\\vspace{0.5pt}
$\bullet$ time dilation time goes slower in moving objects\\
$\bullet$ Lorentz contraction Moving object are shortened only in the direction of travel\\
$\bullet$ 12.1.3 the Lorentz transformations 
Galilean transformations In classical mechanics, events are simultaneou...\\
$\bullet$ Lorentz transformations When we take c as the maximum velocity, we get the transformations:\\
$\bullet$ then the Lorentz-transformations read:\\
$\bullet$ with Λ the Lorentz transformation matrix 
4D scalar product Define a 4-vector as any set of four components...\\
$\bullet$ 4D scalar product Define a 4-vector as any set of four components that transform same as (x 0,x1,x2,x3) und...\\
$\bullet$ > its invariant under Lorentz transformations 
covariant vector aµ = same as the contravariant aµ only with...\\
$\bullet$ > its invariant under Lorentz transformations 
covariant vector aµ = same as the contravariant aµ only with...\\
$\bullet$ - displacement is lightlike: slope is 1 
hyperbola and Lorentz  
transformations\\
\columnbreak


\end{multicols*}

\vspace*{-3pt}
\hrule
\vspace*{2pt}

% PAGE 3: KEY FORMULAS & PROOFS SUMMARY
\begin{multicols*}{3}

\textbf{KEY MAXWELL EQUATIONS:}\\
$\nabla \cdot \E = \rho/\epsilon_0$\\
$\nabla \cdot \B = 0$\\
$\nabla \times \E = -\partial_t \B$\\
$\nabla \times \B = \mu_0 \J + \mu_0\epsilon_0 \partial_t \E$\\

\textbf{DEFINITIONS:}\\
$\D = \epsilon_0 \E + \mathbf{P}$\\
$\H = \B/\mu_0 - \M$\\
$\mathbf{S} = \E \times \B / \mu_0$ (Poynting)\\
$u = (\epsilon_0 E^2 + B^2/\mu_0)/2$ (energy)\\

\columnbreak

\textbf{POTENTIALS:}\\
$\E = -\nabla V - \partial_t \A$\\
$\B = \nabla \times \A$\\
Lorenz gauge: $\nabla \cdot \A + \frac{1}{c^2}\partial_t V = 0$\\

\textbf{BOUNDARY CONDITIONS:}\\
$\hat{n} \cdot (\E_2 - \E_1) = \sigma/\epsilon_0$\\
$\hat{n} \times (\E_2 - \E_1) = 0$\\
$\hat{n} \cdot (\B_2 - \B_1) = 0$\\
$\hat{n} \times (\H_2 - \H_1) = \K_f$\\

\columnbreak

\textbf{RELATIVITY:}\\
$\gamma = (1-v^2/c^2)^{-1/2}$\\
Time dilation: $\Delta t = \gamma \Delta \tau$\\
Length contraction: $L = L_0/\gamma$\\
Velocity add: $v = (u+v')/(1+uv'/c^2)$\\

\textbf{RADIATION (Far-field):}\\
$\E_{rad} \propto \ddot{\mathbf{p}}/r$\\
Larmor: $P \propto a^2$\\
Dipole moment: $\mathbf{p} = q\mathbf{d}$\\

\end{multicols*}

\end{document}
